\section{Differentiated packet calculus and regular ordinary charts}
\label{sec:calculus}

The separation theorem is an assertion about a full tree of successive
expansions.  Its application to equations requires a calculus of those trees,
including a calculus after solving regular equations in the bounded
variables.  We establish that calculus here.  In particular, no simultaneous
convergence of the full exponential series will be used.

The need to verify closure separately from asymptotic injectivity is
substantive: Yeung's analysis of a coefficient class in the classical
finiteness argument exhibits a failure at a differential closure step
\citep{Yeung2025}. We therefore specify the independent arguments and
all differentiated estimates before making an ordinary implicit
substitution. The ensuing closure theorem applies to the packet
hypotheses stated here.

\subsection{The differentiated primitive contract}

Throughout this section a flag, its lifted logarithms, and its positive and
negative angular determinations are fixed.  Write $S$ for the tuple of free
nodes and $t$ for the independent ordinary variables.  A derivative in an
ordinary variable always holds $S$ fixed.  A derivative in a retained free
variable holds the other independent variables fixed; dependent nodes are
then differentiated by their triangular equations.  For a complex input $z$
write
\[
 D_z=\partial_{\Re z},\qquad
 \mathcal A_z=\partial_{\Im z}-i\partial_{\Re z}.
\]
Thus $\mathcal A_z$ is zero on a holomorphic function.  The two lateral
determinations are differentiated separately.

A \emph{finite request} specifies finitely many packet labels, a finite
derivative order, finitely many transition budgets, and a finite Taylor
accuracy.  Its constants, ordinary neighborhoods, angular rooms, and starting
points may depend on the request.  On a selected path every further fixed
request must be available sufficiently far towards infinity on that same
path.  This does not require a common starting point for all requests.
Ordinary tests are real-symmetric holomorphic paths, return to the fixed
ordinary germ at infinity, and have the required positive-order real jets
tending to zero.  Fixed backsteps and slightly larger argument ranges are
available.  The following definition specifies the uniformity needed when
the test itself changes with the target.

\begin{definition}[Controlled families for one finite request]
\label{def:calculus-controlled-family}
Fix a finite request, a compact ordinary box contained with positive
distance in the argument domains of its finitely many labels, and fixed
angular rooms and backsteps.  A controlled family has common bounds for
the finitely many real jets and Taylor-error constants required by this
request, throughout each member's full forward domain beginning at the
prescribed backstep.  The finitely many flag inequalities used there have
common constants and common moduli for their vanishing relative errors.
For convergence near moving anchors, the positive-order ordinary jets
also have a common modulus tending to zero on every fixed window about
those anchors.

By \emph{uniform packet estimates on this controlled family} we mean
that raw growth, side and safe improvements, transitions,
and differentiated source estimates have common constants.  This includes
the number of source terms, their positive precision constants, and the
thresholds at which their specified losses are absorbed.  Uniform
safe-column estimates mean that each member has a suitable column
in the prescribed fixed-length interval; the column itself may depend
on the member.  A common far threshold means thresholds in the finitely
many numerical rate and gap
inequalities used in the request; it need not be a threshold in the
smallest rate alone.  The estimates hold on each full forward domain
once those inequalities and its stipulated ordinary range hold there.
The exceptional constants $D_i$ remain explicit allowed losses.

Every member still returns individually to the fixed germ.  A further
request may require a smaller ordinary box and a different finite set of
controls, reached at a member-dependent later time.  Neither a common
return time to all smaller boxes nor uniformity over infinitely many
requests is part of this definition.
\end{definition}

The common constants in Definition~\ref{def:calculus-controlled-family}
are a quantitative requirement on primitive estimates.  They do not
follow from eventual validity on each separate test.  Below we prove that
separation and the calculus preserve this requirement, and verify that
the interpolation tests belong to such a family.  The primitive estimates
in Sections~\ref{sec:additive}--\ref{sec:mixed} supply this requirement
from norms on common argument boxes.  After those constructions,
Proposition~\ref{prop:calculus-library-uniformity} collects their
quantitative bounds.

\begin{assumption}[Differentiated primitive data]
\label{ass:calculus-primitive}
The primitive packet trees satisfy Definition~\ref{def:packet} and the
following additional requirements.
\begin{enumerate}
\item The initial backgrounds are ordinary coordinates or regular analytic
functions of them.  A raw label in band $j$ is a local smooth function of
ordinary inputs and free inputs whose indices are at least $j$.  Its local
formula, including lifted branches and any choice of cutoff, is the same on
overlapping tests with the same retained determinations.  It has no local
dependence on a discarded free input.  Band~1 functions are holomorphic.
\item Raw growth, its improvement on the band's own side and on safe
columns and their stipulated entering collars, and transitions with
arbitrarily deep owning-fringe errors hold after any fixed number of
independent derivatives.  Transitions are differentiated on both sides,
including their extracted exponential factors.  Every fixed derivative of
$\mathcal A_z$ applied to a label or a background has the source bounds in
Definition~\ref{def:packet}, with a possibly smaller positive constant in the
negative exponent.  These bounds hold locally up to the edges where both
formulas are used.  Pole losses satisfy the same growth and source bounds.
\item The preceding assertions also hold on controlled bounded ordinary
input paths: their real jets are bounded, their required complex values
agree with the corresponding real Taylor continuations to sufficiently
high fixed order, and all argument ranges have fixed spare room.  Slow real
inputs may have their positive-order jets tending to zero.  Safe columns and
the required entering collars can be chosen simultaneously for every finite
collection of labels and derivatives on these paths.  Smooth changes of
cutoff are made locally in the retained variables, consistently on overlaps;
their derivative losses have logarithm $o(U)$ when their precision is
$e^{-cU}$.  For a fixed finite request these estimates hold with the
controlled-family uniformity of
Definition~\ref{def:calculus-controlled-family}, on the full forward
domains, not only at individually selected target points.
These are bounds for the primitive and its independent derivatives
evaluated at the controlled inputs.  A composed path is used in a source
calculation only when its inputs are holomorphic or already have the
differentiated source defects of the previous chart depth.  Finite Taylor
matching by itself gives no exponential bound on a path's anti-CR defect.
\item Supports are iterated left-finite.  More explicitly, the possible
first coordinates form a set finite below each finite ceiling; at each
fixed prefix the possible next coordinates have the same property.  Each
fixed label and its required derivatives extend to every sufficiently
vertical preceding good contour.  There is room for finitely many nested
polynomial angular buffers and for an increase of the side parameter.
\end{enumerate}
\end{assumption}

These are hypotheses on actual primitive evaluations, not on a formal
series alone.  In particular they include safe columns for controlled real
detunings after an ordinary chart substitution.

On a current chart let $b_j$ denote the entire background vector in band
$j$, and put
\begin{equation}
 Z_{k,j}=Z_k(S,b_j),\qquad W_j=e^{-Z_{j,j}}.
 \label{eq:calculus-current-symbol}
\end{equation}
The background invariant is a differentiated transition
\begin{equation}
 b_j=b_{j+1}+\sum_{a>0}W_j^a b_{j,a},
 \label{eq:calculus-background-transition}
\end{equation}
where coefficients belong to the next band.  An equality of this kind means
finite truncations with arbitrarily deep errors on the owning fringe.
The background has controlled real jets and complex Taylor realizations;
at a selected join its difference from the next background is
$O(e^{-cV_j})$.  On the whole owning fringe the difference is
$O(e^{-a_0\Re Z_{j,j}})$ for some $a_0>0$, after narrowing.  Terminal
backgrounds are analytic ordinary germs.

In Equation~\eqref{eq:calculus-background-transition} the extracted factor
$W_j$ is evaluated at $b_j$.  It is never silently replaced by its value at
$b_{j+1}$.  For example, if
\[
 Q(S,x)=\exp(\exp(S+x)),\qquad x=e^{-S-B},
\]
then
\[
 \log\frac{Q(S,x)}{Q(S,0)}
 =e^S(e^x-1)\longrightarrow e^{-B}.
\]
Smallness of $x$ alone therefore does not permit freezing the background of
the extracted $Q$-shift.  Only the rates of indices strictly later than the
current transition are translated below.

\begin{definition}
A packet expression is \emph{positive in series order} if, in each lateral
formal array, every coefficient with lexicographically negative exponent is
the zero analytic germ.  A \emph{strict positive increment} has a
lexicographically positive exponent.  These terms refer to exponent order,
not to the numerical sign of a real function.  The exponent-zero ordinary
germ of a positive expression is denoted by $K_0$.
\end{definition}

\begin{theorem}[Packet calculus]
\label{thm:calculus}
Assume the flag estimates of Section~\ref{sec:flags}, the packet separation
theorem, and Assumption~\ref{ass:calculus-primitive}.  Start with the primitive
data and the initial ordinary backgrounds.  The following operations can be
iterated a finite number of times.
\begin{enumerate}
\item Raw finite sums, products, and any fixed independent derivatives have
differentiated packet trees with iterated left-finite supports.  So does
multiplication by any fixed real exponential monomial
\[
 h^\gamma=\exp\left(-\sum_i\gamma_i Z_i\right).
\]
\item Positive expressions admit cleaned zeroth-prefix realizations $K_j$.
Every fixed ordinary derivative is bounded and converges to the corresponding
derivative of $K_0$ throughout the usable bands.  If a mixed independent
derivative has earliest differentiated free index $l$, then for $j\le l$
its value on $K_j$, multiplied by $e^{C f_l(u)}$, tends to zero for every
fixed $C$.  Differentiated owning transitions satisfy
$K_j-K_{j+1}=O(e^{-c\Re Z_{j,j}})$.  Real local uniform boundedness of a raw
expression implies positivity; its real limit, when specified, identifies
$K_0$.
\item A regular analytic operation on finitely many positive expressions,
whose limiting tuple belongs to its analytic domain, is positive and has
the corresponding differentiated tree.  This includes inversion of a unit
whose limiting value is nonzero.
\item Let $K(S,t,x)$ be a real square vector of positive expressions on a
previous chart.  If
\[
 K_0(t_c,x_c)=0,\qquad \det D_xK_0(t_c,x_c)\ne0,
\]
then the actual regular branch $x=\chi(S,t)$ near $(t_c,x_c)$ and every old
raw expression pulled back to that branch have differentiated packet trees.
The new backgrounds obey Equation~\eqref{eq:calculus-background-transition}.
The ordinary limiting root is analytic.  Band realizations are smooth
regular sheets, holomorphic in band~1, with controlled jets, differentiated
source bounds, safe-column estimates, and differentiated seam bounds.
They retain local independence from discarded free inputs.
\end{enumerate}
For each fixed finite stack of ordinary implicit charts there is a finite
coherence jet order, depending on the stack and fixed angular buffers but
independent of subsequent derivative orders and transition budgets.
Tests on the same free path with the same complex ordinary target and enough
common real jets give the same local sheet.  Every finite request can then
use a larger working Taylor order and farther or narrower working domains;
on overlaps these constructions agree.  All assertions for a fixed finite
request are uniform on controlled families in the sense of
Definition~\ref{def:calculus-controlled-family}.

All formal operations are performed separately in the two signs, by finite
coefficient rules at successive prefixes.  They commute with differentiation,
regular analytic substitution, and the stated regular implicit solution.
Identically zero terminal coefficient germs remain zero under these rules.
\end{theorem}

The proof occupies the rest of this section.  Notice that the theorem makes
a local assertion on lifted determinations.  It asserts no global
single-valuedness over different free-path histories.

\subsection{Support algebra and the current shifts}

\begin{lemma}[Finite coefficient rules]
\label{lem:calculus-support}
Finite unions and Minkowski sums of iterated left-finite subsets of
$\RR^m$ are iterated left-finite.  Each coefficient in their convolution
is a finite sum.  If $A\subset\RR^m$ is iterated left-finite and consists of
strict positives, the semigroup of finite sums of elements of $A$ is
iterated left-finite.  A fixed exponent has only finitely many ordered
representations using nonzero factors from $A$.
\end{lemma}

\begin{proof}
For two supports $A,B$, their first-coordinate sets have lower bounds
$a_*,b_*$.  In a sum with first coordinate at most $M$, the first coordinate
from $A$ is at most $M-b_*$ and the one from $B$ is at most $M-a_*$.  There
are finitely many such choices.  After a first-coordinate sum is fixed,
apply the same argument to each of these finitely many pairs of fibers in
dimension $m-1$.  Induction proves both assertions about convolution.

For the semigroup, induction on $m$ is again appropriate.  A positive
first-coordinate factor has first coordinate at least some $\delta>0$,
since the set of such coordinates is left-finite.  With first-coordinate
budget $M$ there are at most $\lfloor M/\delta\rfloor$ such factors, and
their first coordinates have finitely many possibilities.  Fix one of
these finite patterns.  By the preceding convolution argument their tails
form an iterated left-finite set, bounded below at the next coordinate.
All remaining factors have first coordinate zero and strictly positive
tails.  By the induction hypothesis their tail semigroup is iterated
left-finite and has finite ordered multiplicities at any fixed tail.
Convolve it with the bounded-length tail support already obtained.
There are finitely many interleavings of two fixed finite lists.
This gives both left-finiteness and finite ordered multiplicities.
The empty sum represents zero, and zero factors have been excluded, so it
does not acquire infinitely many repetitions.
\end{proof}

Consequently an analytic Taylor expansion in strict positive increments is
defined coefficient by coefficient.  Its constant term is evaluated as an
ordinary analytic operation; it is not repeated as a zero factor in the
semigroup argument.  At one transition a positive local seed has a smallest
positive exponent.  A fixed owning-coordinate budget therefore requires
only finitely many seed factors and finitely many Taylor derivatives.
Iterated left-finiteness does not assert finiteness of lexicographic initial
segments: for instance $\{(0,n):n\ge1\}\cup\{(1,0)\}$ has an infinite
initial segment below $(1,0)$.  Every finite rule above is taken at successive
fixed prefixes, as the packet construction requires.

Here is the numerical estimate behind these formal operations.  Let a
transition for $F$ and one for $G$ be given on a common $j$-fringe, with
remainders $R_F,R_G$.  A fixed raw label on a sufficiently narrowed fringe
obeys $|F|+|G|\le e^{C\Re Z_{j,j}}$.  Thus
\[
 |R_FG|\le e^{-(M_F-C)\Re Z_{j,j}},\qquad
 |R_GR_F|\le e^{-(M_G+M_F)\Re Z_{j,j}}.
\]
Choose the finite truncations deep enough for the requested budget and the
lower bounds of the supports.  Only after the finite list of resulting
outer labels is known, narrow the fringe so that each of their logarithmic
losses $C'\Re Z_{j+1,j+1}+C'D_j$ is at most
$\eta\Re Z_{j,j}$, with any prescribed fixed $\eta>0$.
This is allowed by the owning-fringe geometry and the path separation.
There is thus no iteration in which a newly generated outer-label constant
forces an endlessly deeper first-coordinate truncation.  The product rule
gives the same estimate for a fixed number of derivatives, after taking
extra depth.  Sums are simpler.  Their sources are still of the permitted
kind because multiplying an $e^{-cU}$ error by factors of logarithmic size
$o(U)$ only decreases $c$.
The distinction between cutoff errors and current-join errors in these
products is recorded in Lemma~\ref{lem:calculus-source-pullback} below.

We next translate later rates.  At a free node $k>j$ the value does not
change with the background.  At a dependent node its exact logarithmic
relation gives
\begin{equation}
 Z_{k,j}-Z_{k,j+1}
 =Z_{k,j+1}\left[
 \exp\left(\sum_{l>k}a_{kl}(Z_{l,j}-Z_{l,j+1})
                   +b_{k,j}-b_{k,j+1}\right)-1\right].
 \label{eq:calculus-rate-difference}
\end{equation}
Work downwards in $k$.  The bracket has strictly positive local $j$-order;
its value and any fixed derivatives are exponentially small compared with
each fixed product of retained later rates.  Taylor expansion of the
exponential therefore gives its differentiated packet transition.
The only numerical scale powers generated by this formula are powers of
dependent nodes.  They reduce exactly to later shifts, since
\begin{equation}
 Z_k^p=e^{p b_k}\prod_{l>k}h_l^{-p a_{kl}}.
 \label{eq:calculus-dependent-power}
\end{equation}
Thus the coefficients in this downward induction are already legitimate
outer packet expressions.  Applying Taylor's formula once more yields
\begin{equation}
 e^{-\gamma Z_{k,j}}
 =e^{-\gamma Z_{k,j+1}}
       \bigl(1+\hbox{strict positive local $j$-increments}\bigr),
 \qquad k>j.
 \label{eq:calculus-later-shift}
\end{equation}
Taylor remainders satisfy the preceding product estimates and their
differentiated versions.  At the current index the factor $W_j^\gamma$
is simply extracted.  Induction from the last transition now constructs
the tree of every $h^\gamma$ and every analytic background factor.
The same calculation shows that changes in a retained raw logarithm
appearing in an earlier cutoff are absolutely tiny at a seam; its numerical
retirement has not been postponed.

\begin{lemma}[Differentiation of the symbols]
\label{lem:calculus-differentiation}
Every fixed derivative of a raw packet expression has the tree obtained by
formal differentiation.  The series for $D\log W_i=-DZ_{i,i}$ has zeros
before coordinate $i$ and has nonnegative local $i$-exponents.  Its later
exponents may be negative.  In addition,
$D_{S_l}W_i=0$ if $l<i$.  Formal differentiation preserves positivity and
cannot move a strict positive exponent to a negative or zero exponent.
\end{lemma}

\begin{proof}
For a free $Z_i$, its independent derivative is a constant or zero.  For a
dependent $Z_i$, triangular differentiation expresses its derivatives as
finite sums of products of dependent powers, background derivatives, and
lower derivatives.  Equation~\eqref{eq:calculus-dependent-power} handles
the powers.  There is a possible apparent recursion: differentiating a
background transition differentiates the extracted $W_j$ as well.
Resolve it first by increasing total derivative order $d$, and at a fixed
$d$ by descending from the later bands.  The terms $D^d Z_{j,j}$ are linear
in the highest derivatives $D^d b_j$, with coefficients made out of lower
derivatives and dependent powers.  In the $d$-fold differentiated form of
Equation~\eqref{eq:calculus-background-transition}, these unknown highest
derivatives occur on the right only after differentiation of a factor
$W_j^a$ with $a>0$.  The resulting finite system has the form
\[
 (I-E_j)D^d b_j=H_j,
\]
where $E_j$ has strictly positive local order and $H_j$ is already known at
this stage.  On a sufficiently far narrowed fringe $\|E_j\|<1/2$.
The inverse $(I-E_j)^{-1}=\sum_{n\ge0}E_j^n$ is a numerical bounded unit
and a formal series with finite coefficients by
Lemma~\ref{lem:calculus-support}.  Truncating it gives the required
differentiated remainder, since the original background transition was
differentiated before being truncated.  This proves the derivative-symbol
rule without assuming it in the background recursion.

For a full monomial the formal rule is
\begin{equation}
 D(c_\lambda h^\lambda)
 =(Dc_\lambda)h^\lambda
   +c_\lambda h^\lambda\sum_i\lambda_i D\log W_i,
 \label{eq:calculus-formal-derivative}
\end{equation}
read successively at the owning transitions.  If $p$ is the first nonzero
coordinate of $\lambda>0$, terms with $i<p$ have multiplier $\lambda_i=0$;
terms with $i>p$ do not alter coordinate $p$; and terms with $i=p$ add a
nonnegative coordinate there.  Hence positivity is preserved.  A zero
coefficient is an analytic germ identity, so its ordinary derivatives and
all its contributions to Equation~\eqref{eq:calculus-formal-derivative}
vanish.  Differentiating each finite actual transition gives exactly these
rules; products and the extra transition depth already established control
the remainders.  The claim follows for higher derivatives by induction.
\end{proof}

\subsection{Negligible defects and cleanup of positive expressions}

\begin{lemma}[Propagation of differentiated source errors]
\label{lem:calculus-source-pullback}
Fix a finite derivative and composition request.  Suppose that each
primitive defect and each input-sheet defect in its chain rules has one
of the permitted differentiated source bounds, and that every remaining
factor has logarithmic size $o(U)$ at the corresponding cutoff source
of cost $U$.  Then the compositions have the same source bounds, with
smaller positive precision constants.  Holomorphic inputs contribute no
input-sheet defect.

At a level-$j$ good join, include all fixed current-level losses in
$C\Re Z_j$, where $\Re Z_j\asymp\eps_jV_j$.  For a transition remainder
$e^{-M\Re Z_j}$, choose the finite budget $M>C$ with the required output
margin before fixing the contour.  The product then has precision
$e^{-(M-C)\Re Z_j}$, hence $e^{-cV_j}$ for a positive fixed $c$ on that
contour.  For an independently established error $e^{-c_0V_j}$ with
$c_0>0$ fixed independently of the new contour choice, choose $\eps_j$
sufficiently small that $C\Re Z_j\le c_0V_j/2$.
In either case the remaining later-rate and fixed derivative losses of
logarithmic size $o(V_j)$ can be absorbed afterwards.  For a controlled
family, assume that the defect bounds and the $o(U)$ and $o(V_j)$ controls
are common for this finite request.  Then all these choices are common
to the family.
\end{lemma}

\begin{proof}
For an input map $\Psi=(\Psi_a)$, Wirtinger differentiation in any one
source coordinate gives
\[
 \bar\partial(F\circ\Psi)
 =\sum_a(\partial_aF\circ\Psi)\bar\partial\Psi_a
  +\sum_a(\partial_{\bar a}F\circ\Psi)
                         \overline{\partial\Psi_a}.
\]
Here the derivatives of $F$ are in its independent input coordinates.
Thus every term contains an input-sheet defect or an ambient primitive
defect.
After any fixed further derivatives there are finitely many such terms,
and each still contains a differentiated defect.  Its other factors
multiply $e^{-cU}$ by $\exp(o(U))$, which decreases the positive constant
$c$ but preserves the allowed precision.  The two stated join estimates
are the corresponding product estimates with the current loss kept
explicit.  A fixed $\eps_jV_j$ need not be $o(V_j)$, so reducing the
collar alone cannot compensate for a transition budget $M\le C$.
Finite Taylor agreement of an input supplies geometric control only;
its differentiated source bounds must already be available for this
chain-rule argument.
\end{proof}

We record how the differentiated source errors are used at a transition.
For every finite $M$, after the permitted narrowing and increase of the
side parameter,
\begin{equation}
 e^{-cU}\,\exp(o(U))
       =O(e^{-M\Re Z_{j,j}})
 \label{eq:calculus-source-fringe}
\end{equation}
for each source active on either side of the $j$-transition, including any
fixed path-derivative losses.  Here is the comparison of scales.  For a
source from the next band with $f_r\le f_j^{3/2}$, the angular estimate of
Lemma~\ref{lem:flags-path} gives an absolute $o(1)$ loss in the logarithm
of the cost.  If a secondary index $k>j$ occurs, its cosine is bounded
below by a constant times $(\Lambda_j-\Lambda_k)/\Lambda_j$; the gap
estimates pay its logarithm.  More explicitly, on a fringe with
$\cos(\Im f_j)\le\eps/\Lambda_j$, the simple and mixed costs respectively
satisfy
\[
 \frac{U}{\Re Z_j}\ge
 c e^{f_r-f_j}\frac{\Lambda_j}{\eps},\qquad
 \frac{U}{\Re Z_j}\ge
 c e^{f_r-f_j}\frac{\Lambda_j-\Lambda_k}{\eps}.
\]
The gap derivative tends to infinity; these quotients therefore tend to
infinity even when $r=j$.  If $f_r>f_j^{3/2}$, the logarithm of the cost
dominates $f_j$ even with this loss.  On the inner side, a cost with
$\sigma(r)=j$ is made sufficiently strong by increasing the side
parameter.  When its secondary index is $j$, the quotient of its scale by
$\Re Z_{j,j}$ contains the diverging factor $e^{f_r-f_j}$ with the same
cosine on numerator and denominator.  Earlier leading indices give still
greater suppression.  These are precisely the allowed source cases.
Finally, a fixed derivative of a retained free path of index $l\ge j$
is bounded by $e^{C f_l(u)}$ by Lemma~\ref{lem:flags-path}; these losses
have logarithm $o(U)$.  This proves
Equation~\eqref{eq:calculus-source-fringe}.  In particular a finite smooth
Taylor computation may discard its anti-CR terms at any prescribed
owning-coordinate precision, but must first bound those terms as above.

\begin{lemma}[Differentiated cleanup]
\label{lem:calculus-cleanup}
Let $K$ be positive.  Keep its raw zero-prefix packets
$K_j=S^j_{0,\ldots,0}$, inserting a zero packet if necessary, and in each
of their transitions delete all negative local exponents.  The resulting
transitions still hold with arbitrarily deep errors and all requested
independent derivatives.  Their zero-prefix derivative packets are exactly
the derivatives of $K_j$.
\end{lemma}

\begin{proof}
Fix a finite transition and derivative request.  A negative prefix through
level $j$ has an entirely zero formal subtree, and every earlier subtree
is also zero.  The zero-subtree conclusion of
Theorem~\ref{thm:separation} therefore bounds its raw outer coefficient by
$O(e^{-cV_j})$ on a sufficiently vertical $j$-collar and the corresponding
fringe.  Lemma~\ref{lem:calculus-differentiation} says the same about the
formal derivative tree of $K$.

To obtain a bound for actual derivatives of an individual deleted
coefficient, retain the raw labels before canceling zero terminal leaves
and use induction on derivative order.  In the formal expansion
of a derivative of $K$, the contribution at its negative prefix contains
that derivative of the coefficient itself.  Every other contribution has
fewer derivatives on a coefficient at a negative prefix; its remaining
factors arise from differentiated shifts.  Positive prefixes cannot
contribute, by Lemma~\ref{lem:calculus-differentiation}, and the all-zero
prefix has no shift derivative.  There are finitely many relevant terms
at the fixed budget.  The induction hypothesis gives $e^{-cV_j}$ for their
coefficient derivatives, whereas the remaining factors have logarithmic
size $o(V_j)$ on this fringe.  Subtracting these terms from the derivative
packet proves the assertion for the next derivative order.  The anti-CR
parts are covered by Equation~\eqref{eq:calculus-source-fringe}.

For the finitely many deleted local exponents $a<0$, choose the good
contour sufficiently vertical that
$|a|\Re Z_{j,j}$, the requested $M\Re Z_{j,j}$, and all fixed losses are
smaller than, say, $cV_j/2$.  Multiplying the deleted coefficients by
$W_j^a$ then leaves an $O(e^{-M\Re Z_{j,j}})$ error with the requested
derivatives.  The coupled leading estimates and the original transition
extend this comparison through the usable owning fringe.  This proves the
cleaned transition.  A strict positive earlier exponent cannot feed into a
zero prefix on differentiation, so the zero-prefix derivative packet is
the actual derivative of $K_j$, as claimed.
\end{proof}

\begin{lemma}[Uniform separation for a controlled family]
\label{lem:calculus-uniform-separation}
For one controlled family and one fixed finite coupled calculation,
the coupled growth and zero-subtree upgrade constants of
Section~\ref{sec:separation} are common to the family.  Leading limits
are uniform at its moving targets.  Deeper preliminary estimates used to
prove a zero-subtree upgrade may have member-dependent constants and
starting points, including dependence on the return time to the germ.
\end{lemma}

\begin{proof}
Definition~\ref{def:calculus-controlled-family} supplies precisely the
common finite-label, source, and flag controls in
Lemma~\ref{lem:separation-family-uniformity}.  Apply that lemma to the
fixed coupled calculation and its finitely many derivative packet trees.
Its weighted estimate removes the individual preliminary constants and
return times; its terminal comparison uses the common slow-variation
modulus near the moving targets.  This gives the stated common constants
and uniform convergence at moving targets.
\end{proof}

\begin{lemma}[Bounds, limits, and free derivatives]
\label{lem:calculus-positive-bounds}
The cleaned realizations of a positive expression satisfy all the bounds in
part~(2) of Theorem~\ref{thm:calculus}.  Conversely, real locally uniform
boundedness of a raw expression on an anchored ordinary neighborhood
implies positivity.
\end{lemma}

\begin{proof}
Apply the coupled leading conclusion of Theorem~\ref{thm:separation} with
zero prefix to the cleaned expression and its ordinary derivatives.
If its exponent-zero coefficient is absent, insert it as zero.  This yields
boundedness and convergence to the terminal exponent-zero germ $K_0$
inside the chosen contours; the cleaned transition gives the same result
to the polynomial side after increasing its side parameter.  There is a strict gap to the
next positive local exponent, so that transition and all its fixed
derivatives give
\[
 K_j-K_{j+1}=O(e^{-c\Re Z_{j,j}}).
\]
Equality of the two exponent-zero germs follows from separation.
All these statements are uniform on the controlled families of
Definition~\ref{def:calculus-controlled-family}, by
Lemma~\ref{lem:calculus-uniform-separation}.  In particular this controls
families of different tests at different anchors; pathwise eventuality
alone is not being used for this step.

Now differentiate in a free input and let $l$ be the earliest index of any
free input differentiated.  A term whose first active index is after $l$
has no dependence on $S_l$, because its coefficients use only retained
inputs.  A terminal coefficient also has no free dependence.  Thus every
nonzero formal derivative term has a strict positive first index at or
before $l$.  In the cleaned zero-prefix packet at level $j\le l$, the
first such index lies between $j$ and $l$.  Normalize by its leading
positive shift and use separation.  On the relevant domain
$\Re Z_k/f_l\to\infty$ for $j\le k\le l$, by
Equation~\eqref{eq:flags-retained-buffer} in
Lemma~\ref{lem:flags-contours}.  Hence that shift absorbs
$e^{C f_l}$ for every fixed $C$.  If there is no nonzero term in this
interval, request an arbitrary positive dummy exponent at level $l$ in the
zero-subtree estimate.  Positive exponents at earlier levels are later in
lexicographic order than this dummy request and cause no obstruction.
This proves the asserted decay, also after further ordinary and free
derivatives.  A product of finitely many free-path derivative factors is
at most $e^{C'f_l}$, so the same estimate pays all chain-rule losses.

For the converse, if a negative coefficient is nonzero, iterated
left-finiteness supplies a first nonzero negative exponent $\lambda$.
Choose an ordinary point at which its nonzero analytic coefficient is
nonzero.  Separation gives
$K=h^\lambda(c_\lambda+o(1))$ there.  The first nonzero coordinate of
$\lambda$ is negative, and the ordered real scales imply
$|h^\lambda|\to\infty$.  This contradicts local uniform boundedness.
Once positivity is known, its real limiting germ is its zero coefficient
by the first part of the proof.
\end{proof}

This converse is useful after a normalization: multiplication by a monomial
is first a raw operation, and its positivity can then be checked by actual
real bounds.  No positivity assumption on the expression before division
is required.

\begin{lemma}[Regular analytic operations]
\label{lem:calculus-analytic}
Let $H$ be analytic near the limiting tuple of finitely many positive
expressions.  Applying $H$ to their actual values produces a positive
expression with the formal analytic substitution and the differentiated
packet estimates.  This includes a matrix inverse at an invertible
limiting matrix.
\end{lemma}

\begin{proof}
In a band the cleaned tuple stays in a fixed smaller neighborhood of the
limiting tuple by Lemma~\ref{lem:calculus-positive-bounds}.  At a transition
write it as the outer tuple plus a strict positive increment $\delta$.
Taylor's formula with integral remainder gives
\[
 H(y+\delta)=\sum_{|\alpha|<N}
       \frac{\partial^\alpha H(y)}{\alpha!}\delta^\alpha
       +O(\|\delta\|^N).
\]
All derivatives of $H$ required by a finite request are bounded on the
smaller neighborhood.  Since $\delta$ and its fixed derivatives are
exponentially small on the owning fringe, take $N$ large enough and apply
the product estimates to control the differentiated remainder.  The
coefficients are outer packet operations, to which the same procedure
applies recursively.  Lemma~\ref{lem:calculus-support} gives the formal
finite coefficient rule.  Source errors retain their precision under
bounded derivatives of $H$, and finite intersections of the primitive safe
sets suffice.  The actual result is bounded on the real neighborhood and
therefore positive by Lemma~\ref{lem:calculus-positive-bounds}.
\end{proof}

\subsection{The real implicit branch and its trial inputs}

It remains to prove regular implicit closure.  A dependent exponential
may rotate many times on a fixed complex disk in an ordinary variable,
so we cannot assume that such a disk is an admissible implicit domain.
We will first solve the real equation, then use its Taylor jets to
construct complex roots on smaller domains.

Proceed by induction on the number $d$ of regular ordinary charts already
constructed.  At depth $d$ assume there is an integer $r_d$, fixed by the
stack and its angular rooms, independently of later derivative orders
and transition budgets, such that tests on the same free path, with a common complex ordinary target
and common real input jets through $r_d$, give the same local stack of
solved sheets.  On each test the solved variables agree with their real
Taylor jets to every prescribed finite order on sufficiently far working
subdomains.  They have all the required fixed path and independent
derivative bounds, differentiated sources, and exponentially small seams.
These assertions are uniform for every fixed finite request on controlled
families.  At depth zero they are the primitive local-determination and
controlled-path hypotheses.

Consider $K(S,t,x)=0$ as in Theorem~\ref{thm:calculus}, and put
$J_0=D_xK_0(t_c,x_c)$.  Shrink the ordinary neighborhood so that
$D_xK_0(t,x)$ stays uniformly close to the invertible matrix $J_0$.  By
Lemma~\ref{lem:calculus-positive-bounds}, the real equation for sufficiently
far real inputs has a unique nearby branch $x=\chi(S,t)$.  For example,
the map $x\mapsto x-J_0^{-1}K(S,t,x)$ contracts a fixed sufficiently
small ball and maps it into itself, because $K$ and $D_xK$ converge to
$K_0$ and $D_xK_0$.  The real analytic implicit-function theorem then gives
the actual branch.  Differentiating its exact equation yields bounded
fixed ordinary derivatives.  Along a slow real test its positive-order
path jets tend to zero: each chain-rule term either contains a small
ordinary path derivative or a free derivative paid by
Lemma~\ref{lem:calculus-positive-bounds}.

To continue this branch in band $j$, we will solve $K_j=0$ near a high
Taylor value of $\chi$.  Before estimating that equation, we need the old
chart at every trial input in a small tube about the Taylor value.
The following interpolation keeps the lower real jets which identify
the old sheet, while attaining each complex trial input exactly.

\begin{lemma}[Interpolation with fixed lower jets]
\label{lem:calculus-interpolation}
Fix an integer $r$ and a finite higher jet order.  Let $v\ne0$ be small.
Suppose a bounded ordinary input has prescribed real jets at $u$ through
that higher order, with its positive-order jets small, and the proposed
complex target differs from the associated high Taylor value at $u+iv$
by $O(|v|^R)$, where $R>r+2$.  After initially matching sufficiently many
jets, the target can be attained exactly by a bounded real-symmetric
holomorphic test that keeps all jets through order $r$.  The test returns
to the prescribed ordinary germ on the real ray.  Its additional
coefficients tend to zero as $v\to0$.
For a family with common finite jet bounds, target-error constants, and
ordinary slack, the interpolants have common bounds for every fixed
requested number of derivatives on a full forward strip.  If the prescribed
positive real jets and $|v|$ tend uniformly to zero, these derivative
bounds have a common vanishing modulus.  Their times of return to a
smaller neighborhood of the germ need not be uniform.
\end{lemma}

\begin{proof}
It suffices to treat one scalar component.  On a fixed narrow strip put
\[
 B(s)=1-\tanh^2(c_0(s-u)),\qquad
 \phi_n(s)=B(s)\tanh^n(c(s-u)),
\]
where $c>0$ is fixed and small, and $c_0>0$ can be made still smaller.
Both functions are holomorphic on a strip with fixed extra room,
real-symmetric, bounded there, and $\phi_n(s)$ tends to zero as
$s\to+\infty$ on the real ray.  Its Taylor coefficient of degree $n$ at
$u$ is $c^n\ne0$, and lower coefficients vanish.  A finite triangular
linear system therefore matches any prescribed finite real jet.  A real
value offset from the germ is supplied by $B$; its positive derivatives
are made small by choosing $c_0$ small.  The coefficients matching the
remaining positive-order jets are small.  Taking the ordinary neighborhood
smaller preserves the argument range, with slack.

Let $n_e,n_o$ be the first even and odd integers greater than $r$; then
$\max(n_e,n_o)\le r+2$.  At $s=u+iv$, $B$ is real and
$\tanh(civ)=i\tan(cv)$.  Thus $\phi_{n_e}(u+iv)$ is real nonzero and
$\phi_{n_o}(u+iv)$ is purely imaginary nonzero.  Two real coefficients
therefore correct independently the real and imaginary target errors,
without changing the prescribed lower jets.  Their sizes are at most
\[
 C|v|^{R-n_e}+C|v|^{R-n_o}
 \le C'|v|^{R-r-2}=o(1).
\]
If required, start with a still higher matched jet so its own Taylor
remainder is covered by this target error.  All corrections decay along
the ray, so the return to the germ is unchanged.  The construction applies
componentwise to vector inputs.

We check the full-domain quantitative assertion.  Fix a strip
$|\Im(s-u)|<\sigma$ strictly inside the poles for the chosen fixed $c$,
and keep $0<c_0\le c$.  For every fixed pair of integers $L,k$, the
functions $B\tanh^n(c(s-u))$, $n\le L$, and their first $k$ derivatives
are bounded there by constants depending only on $L,k,c,\sigma$.
These constants are independent of $u$ and $c_0$.  Derivatives of positive
order of $B$ alone are bounded by $C_k c_0$.  The diagonal entries of the
finite jet matrix are $n!c^n$, independent of $c_0$; its other entries
are bounded.  Its inverse therefore has a common bound for
$0<c_0\le c$.

Write $a_0$ for the value offset from the germ and let
$a_1,\ldots,a_L$ be the prescribed positive-order real jets.  The initial
jet-matching coefficients satisfy
\[
 \max_{1\le n\le L}|b_n|
 \le C_L\left(\max_{1\le n\le L}|a_n|+c_0|a_0|\right).
\]
The two final correcting coefficients have the uniform bound already
proved, since $|\phi_n(u+iv)|\ge c_n|v|^n$ for sufficiently small $|v|$,
uniformly for $0<c_0\le c$.  Consequently the final interpolant $T$ obeys
\begin{equation}
 \sup_{\substack{\Re s\ge u-B_*\\|\Im s|<\sigma}}
       |\partial_s^kT(s)|
 \le C_{L,k}\left(c_0|a_0|
        +\max_{1\le n\le L}|a_n|
        +|v|^{R-r-2}\right),\qquad k\ge1.
 \label{eq:calculus-interpolation-family}
\end{equation}
Here $B_*$ is any fixed backstep, and $L$ is first chosen large enough
for the fixed jet and remainder request.  The value bound is common as
well.  On the real line the value-offset term lies on the segment from
the germ to the prescribed value; on the narrow strip its excess and
the remaining small terms are absorbed by the fixed ordinary slack.
Taylor's integral remainder and the next derivative bound in
Equation~\eqref{eq:calculus-interpolation-family} give common complex
Taylor-error constants.  Choosing $c_0$ with a common modulus tending to
zero proves the claimed vanishing derivative modulus on the entire forward
domain, while every individual bump still decays to zero at infinity.
This proves membership in Definition~\ref{def:calculus-controlled-family}
whenever the free paths satisfy its common flag controls.

No derivative of the interpolation coefficients with respect to the
complex target is used.  Those derivatives can contain powers of
$|v|^{-1}$.  The estimates concern derivatives of the test as a function
of $s$; intrinsic derivatives on the old graph are instead obtained from
its already established packet or implicit equations.
\end{proof}

Band domains can be prescribed before the new complex root is known.
Use the actual free paths and sufficiently high Taylor jets of the real
backgrounds as reference backgrounds.  By
Lemma~\ref{lem:flags-background}, on $|v|\le C/\Lambda_i(u)$ a perturbation
$\Delta b=O(|v|^q)$ changes a relevant lifted logarithm by at most
\begin{equation}
 |\Delta f_i|
 \le C\Lambda_i\bigl(1+\log(2+\Lambda_i)\bigr)^C |v|^q.
 \label{eq:calculus-angle-stability}
\end{equation}
The estimate holds throughout the short perturbation segment.  Its
triangular proof uses that the derivative contributions from the children
of a dependent node are dominated by its first child
$Z_{\sigma(i)}\Lambda_{\sigma(i)}$; a free node has no background
variation.  Choosing $q$ large makes the right side smaller than the
polynomial angular buffer and the fixed good-collar width.  We may use
several strictly nested rooms.  This makes the reference domains
independent of an unknown implicit continuation.

The interpolation lemma supplies a controlled family of trial inputs,
rather than separate tests with unrelated eventual starting points.
At depth $d$, its fixed lower-jet agreement identifies one old graph
throughout the trial tube.  The construction below preserves this
controlled-family assertion using the finite coefficient rules, uniform
separation, and common regular inverse bounds.

The coherence order identifies the local sheet.  A larger working Taylor
order will control the accuracy of a particular finite calculation.
Increasing that order must leave the identified sheet unchanged.

\subsection{Complex roots and coherence}

First choose a broad comparison-tube order $R$ using only $r_d$,
Equation~\eqref{eq:calculus-angle-stability}, and first-derivative control
of the finitely many old equations.  Choose it with a finite margin,
in particular larger than the jet orders needed for interpolation at the
old depth.  This choice will determine the next coherence threshold.
Only now, for a given finite request, choose an arbitrarily larger working
Taylor order $N$ and put
\begin{equation}
 p_N(v)=\sum_{a=0}^{N-1}\frac{\chi^{(a)}(u)}{a!}(iv)^a,
 \label{eq:calculus-taylor-center}
\end{equation}
where these are real path derivatives along the actual composite real
branch.  The old chart can be evaluated at every point of
$\|x-p_N(v)\|\le C_0|v|^N$ and on the broader comparison tube of order
$R$.  Indeed interpolation gives the required input tests, and old
coherence puts them on the same old graph.  The older solved variables
have common lower Taylor centers since their real equations have the
prescribed common real jets.  The old regular Jacobians give one graph
throughout these nested tubes, including the segments between trial
points.  This statement uses no new root.
For a controlled family the prescribed real solution jets have the common
bounds and vanishing moduli just obtained by real implicit
differentiation.  Equation~\eqref{eq:calculus-interpolation-family}
therefore puts all these trial inputs, for all its anchors and targets,
in one controlled family at the old depth.  The old derivative estimates
and limiting Jacobian comparison apply uniformly to the entire trial
tube.  Only the finite request fixes $N$ and the numerical far thresholds;
the individual interpolating test does not change those choices.

Consider the residual on that old graph,
\[
 R_j(v)=K_j(S(u+iv),t(u+iv),p_N(v)).
\]
Its derivatives in $v$ through order $N$ are bounded on each used band,
with one bound for the controlled family at this fixed request.
The ordinary chain-rule factors are bounded, and every free-path factor
is absorbed by the independent free-derivative estimate of
Lemma~\ref{lem:calculus-positive-bounds}.  In the first band the residual
is holomorphic near zero, and its first $N$ Taylor coefficients vanish
because the real equation and the jets in
Equation~\eqref{eq:calculus-taylor-center} agree.  At a selected join
$v_i$, the residual and all its requested derivatives have exponentially
small jumps, even after the finite path-factor losses.

For clarity, repeated integration on the successive intervals gives the
following elementary piecewise Taylor estimate.  If $R^{(a)}(0)=0$ for
$0\le a<N$, $\|R^{(N)}\|\le C$, and its derivative jumps at $v_i$ are
$\Delta_{i,a}$, then
\begin{equation}
 \|R(v)\|
 \le \frac{C|v|^N}{N!}
   +\sum_{i:\,|v_i|\le |v|}\sum_{a=0}^{N-1}
       \frac{\|\Delta_{i,a}\|\,|v-v_i|^a}{a!}.
 \label{eq:calculus-piecewise-taylor}
\end{equation}
Each seam bound is $o(|v_i|^{N-a})$, because exponential seam precision
beats every fixed power of its angular location.  Since
$|v-v_i|\le |v|$ on a fixed lateral vertical segment, each summand is
$O(|v|^N)$.  There are finitely many joins.  Consequently
\begin{equation}
 R_j(v)=O(|v|^N)
 \label{eq:calculus-residual}
\end{equation}
throughout the band up to its used side.  This argument uses high real
jets and bounded derivatives, not a Cauchy estimate on a fixed ordinary
disk.

Identify $\CC^{\dim x}$ with a real vector space, and write $J_0^{\RR}$
for the real representation of $J_0$.  On the entire broad comparison
tube the ordinary derivative converges uniformly to the corresponding
derivative of $K_0$, and the anti-CR part is uniformly negligible.
The difference from $J_0$ is the sum of this convergence error and
$D_xK_0(t,x)-J_0$; the second is uniformly small on a sufficiently
small fixed ordinary neighborhood.  Consequently a common far threshold
gives
\begin{equation}
 \left\|I-(J_0^{\RR})^{-1}D_x^{\RR}K_j\right\|\le\tfrac12.
 \label{eq:calculus-contraction}
\end{equation}
The map $T_j(x)=x-(J_0^{\RR})^{-1}K_j(x)$ contracts there.  By
Equation~\eqref{eq:calculus-residual}, for a sufficiently large fixed $C_0$
it maps the closed ball of radius $C_0|v|^N$ about $p_N(v)$ into itself.
It has exactly one root $x_j$ in that ball.  Furthermore any two roots
in the broad tube coincide, by integrating
Equation~\eqref{eq:calculus-contraction} on their connecting segment.
The first-band root is holomorphic by the holomorphic implicit-function
theorem.

At any nonzero band point the ordinary smooth implicit-function theorem
gives a local smooth continuation of this root.  Such a neighborhood may
be arbitrarily small.  The real jet centers vary continuously with the
base point; the available lower-order tube has slack.  The local
continuation therefore remains in the same broad tube and agrees with
the tested root there by uniqueness.  This proves that the constructed
objects form smooth local sheets, rather than unrelated pointwise roots.
Near $v=0$ the actual real and first-band implicit branches give the same
conclusion.  At the terminal edge the limiting analytic equation can be
evaluated first at $x_m$, and its ordinary regular inverse yields
$x_{m+1}$; no extension of an old band sheet past that edge is required.

To check coherence, take two tests on the same free path, with the same
complex ordinary target and enough common real input jets.
Their real implicit solution jets agree to the corresponding order,
since successive differentiation of the regular equation determines
those jets uniquely.  Choose this fixed order high enough that both
roots lie in the same broad tube of order $R$.  Old coherence and
interpolation identify the previous sheets there, including on the
segment between the roots.  Equation~\eqref{eq:calculus-contraction}
then identifies the new roots.  This defines a finite $r_{d+1}$ using
only $r_d$, the angular margin, and first-derivative regularity.  A later
construction with larger $N$ lies in the same broad tube sufficiently
far out, so it gives the same root.  Thus branch identification does not
consume the arbitrary derivative accuracy requested afterwards.

\subsection{Differentiated sheets and their transition series}

We give the remaining estimates for the new graph explicitly.  For an
independent real coordinate $y$, its exact first derivative satisfies
\begin{equation}
 D_yx_j=-(D_x^{\RR}K_j)^{-1}D_yK_j.
 \label{eq:calculus-implicit-derivative}
\end{equation}
For a higher multi-index $\alpha$, differentiating the equation gives
\[
 D_x^{\RR}K_j\,D^\alpha x_j
 =-D_y^\alpha K_j-\mathcal P_\alpha,
\]
where $\mathcal P_\alpha$ is a finite sum of products of old derivatives
and lower derivatives of $x_j$.  The inverse is uniformly bounded by
Equation~\eqref{eq:calculus-contraction}.  Induction proves bounded
ordinary derivatives.  If a free derivative with earliest index $l$
occurs, each term contains either an old derivative with such a free
index or a lower graph derivative with that index; induction and
Lemma~\ref{lem:calculus-positive-bounds} give decay sufficient to pay
$e^{C f_l}$ for any fixed $C$.  Applying $\mathcal A_y$ to the exact
equation gives the same invertible real linear system with right side
consisting of old anti-CR defects.  Higher differentiated defects follow
from the same finite induction: each term contains an old differentiated
defect, and its other factors have the established absorbable losses.
Lemma~\ref{lem:calculus-source-pullback} therefore preserves the old source
precision after decreasing its exponent constant.  Replacing the real Jacobian by its
complex-linear part changes these formulas only by errors of that
precision.

At a $j$-join compare the two equations on their common old graph.
Evaluate the outer one at the inner root.  The cleaned transition gives
\[
 K_{j+1}(x_j)
 =K_{j+1}(x_j)-K_j(x_j)
 =O(e^{-c\Re Z_{j,j}(x_j)}).
\]
The segment from $x_j$ to $x_{j+1}$ lies in the comparison tube.  Its
Jacobian is uniformly regular, so
\begin{equation}
 d_j:=x_j-x_{j+1}
   =O(e^{-c\Re Z_{j,j}(x_j)}).
 \label{eq:calculus-root-seam}
\end{equation}
For derivatives, subtract Equation~\eqref{eq:calculus-implicit-derivative}
and its higher versions on the two sides.  Differences of old derivatives
at the same input are bounded by their differentiated transition; their
differences at $x_j$ and $x_{j+1}$ are bounded by the next old derivative
integrated on the segment.  The inverse difference is
$A^{-1}-B^{-1}=A^{-1}(B-A)B^{-1}$.  Induction therefore proves
Equation~\eqref{eq:calculus-root-seam} with every requested independent
derivative, using extra transition depth to pay fixed triangular losses.
This is a differentiated equation argument; it does not differentiate a
pointwise smallness assertion.  The resulting path-derivative seam bounds
and Equation~\eqref{eq:calculus-piecewise-taylor} also give arbitrary
fixed Taylor accuracy for the derivatives of the new sheets.

At band $j$, all the equations on the old graph use only retained free
inputs.  Differentiating in a discarded free direction has zero right
side in Equation~\eqref{eq:calculus-implicit-derivative}; uniqueness
therefore gives zero derivative of the new graph in that direction.
This proves local retained-input independence.  A discrete continuation
choice may have used the free path, but it introduces no local derivative
and is fixed on the lifted germ.

We now construct the transition series, keeping
\[
 W'_j=W_j(S,t,x_j)
\]
as the current extracted symbol.  All new outer coefficients are evaluated
at $x_{j+1}$.  Taylor-expand the old outer coefficients along the segment
from $x_{j+1}$ to $x_j=x_{j+1}+d_j$.  Smooth complex Taylor expansion uses
$D_x$; the terms containing anti-CR derivatives have arbitrarily deep
fringe precision by Equation~\eqref{eq:calculus-source-fringe}.  The
segment is exponentially short.  Its rate and retained cutoff logarithms
are stable by the triangular estimates, so fixed losses are absorbable
relative to $\Re Z_{j,j}(x_j)$.  Let
$J=D_xK_{j+1}(x_{j+1})$.  The finite Taylor equation, to any specified
budget, has the form
\begin{equation}
 0=Jd_j+\sum_{h\ge2}B_{0,h}[d_j^h]
       +\sum_{a>0,\,h\ge0}(W'_j)^a B_{a,h}[d_j^h].
 \label{eq:calculus-formal-implicit}
\end{equation}
Brackets denote the corresponding symmetric multilinear evaluations in
the vector case; both sums are truncated as required, with a remainder
whose independent derivatives have the desired precision.  This follows
by applying the old differentiated transition at $x_j$ and Taylor's
integral remainder on the common segment, using the already proved
differentiated smallness of $d_j$.

Solve Equation~\eqref{eq:calculus-formal-implicit} formally for
\begin{equation}
 d_j=\sum_{a>0}(W'_j)^a d_{j,a}.
 \label{eq:calculus-implicit-series}
\end{equation}
Multiply by $J^{-1}$.  Every term other than the linear term either
contains a positive seed exponent or contains at least two factors of
$d_j$.  The positive seed support has a smallest positive local exponent
and is left-finite.  At a given owning budget, its number of factors and
the depth of the recursion are bounded.  Thus
Lemma~\ref{lem:calculus-support} gives a unique finite coefficient
calculation for each $d_{j,a}$.  Every such coefficient uses finitely many
old outer coefficients and derivatives, and the bounded unit $J^{-1}$.
The inverse itself has an outer transition with the same limiting
inverse and strict positive increments by
Lemma~\ref{lem:calculus-analytic}.

To verify that the formal answer describes the actual graph, truncate
Equation~\eqref{eq:calculus-implicit-series} at a sufficiently deep finite
order and call its value $d_*$.  Crucially, evaluate $d_*$ at the same
current numerical $W'_j$ as the exact equation.  Formal cancellation and
the finite remainder estimates show that $d_*$ satisfies the finite
polynomial equation with an error of the requested differentiated
precision.  The exact $d_j$ satisfies it with the same precision.
Subtracting yields
\[
 (J+E)(d_j-d_*)=R,\qquad
 \|E\|\longrightarrow0,
\]
where $R$ has that precision with all requested independent derivatives.
The bounded inverse and repeated differentiation of this identity give
the same precision for $d_j-d_*$.  Fixed derivatives of $W'_j$ incur only
the previously established triangular losses; taking additional initial
depth pays them.  This proves the actual differentiated transition
series without assuming that pointwise asymptotics can be differentiated.

For an arbitrary old raw packet $F$, apply its old transition at $x_j$
and Taylor-expand its outer coefficients at $x_{j+1}$ using
Equation~\eqref{eq:calculus-implicit-series}.  Its negative exponents only
require finitely many extra orders at any fixed prefix.  The resulting
outer coefficients again use the same finite grammar.  In particular
$b'_j=b_j(S,t,x_j)$ has the strict positive background transition.
Terminal roots solve the common regular analytic equation and therefore
give analytic coefficient germs.

Finally, the new raw growth, sources, and safe estimates follow from
these finite formulas and the strengthened primitive condition.  At the
bottom of a formula its non-unit raw factors are primitive evaluations
on the bounded controlled coordinates of the stack.  The real coordinates
are slow, the complex coordinates have the just-proved high real-jet
control, and finite simultaneous safe columns are available by
Assumption~\ref{ass:calculus-primitive}.  Positive factors and inverses
are uniformly bounded on the comparison tubes by interpolation and the
regular Jacobians.  Fixed derivative losses have already been paid in
the source comparison.  This completes the induction on chart depth and
the proof of Theorem~\ref{thm:calculus}.

\subsection{Formal compatibility and refinement}

For precision we collect the formal conclusion used in the zero argument.
At every finite prefix, addition and multiplication are finite convolution
rules, differentiation is Equation~\eqref{eq:calculus-formal-derivative},
analytic substitution is Taylor expansion in strict increments, and a
regular ordinary implicit substitution is the uniquely solved recursion
in Equation~\eqref{eq:calculus-formal-implicit}.  These descriptions hold
separately in the two signs and coincide with the actual differentiated
transitions.  The ordinary exponent-zero coefficients are evaluated on
the unique regular analytic limiting branch.  Hence they respect analytic germ
identities, including identities obtained by differentiation.  Applying the
same well-defined operations to both sides preserves a formal identity.
In particular differentiation, multiplication, and pullback preserve a
formally zero expression, and Theorem~\ref{thm:separation} identifies the
resulting actual zero.  This statement concerns the full packet tree; it does not
identify unrelated numerical realizations by discarding a flat error.

\begin{lemma}[Retesting a finite chart recipe]
\label{lem:calculus-refinement}
Enlarge a saturated flag by finitely many permitted insertions without
converting its old active free nodes.  Suppose the primitive data can be
retested on this flag with Assumption~\ref{ass:calculus-primitive}, the same
terminal analytic germ recipes and lateral determinations, and identity
transitions at unused new nodes.  Then any previously constructed finite
ordinary chart recipe can be retested on the refinement.  Its real
expressions on the old anchored inputs are unchanged, and its positivity,
normalizing divisions, regular implicit steps, and formal compatibility
remain valid.
\end{lemma}

\begin{proof}
Retest the old primitives before making a new ordinary substitution.
Initially the bounded parameters introduced by the refinement are passive
ordinary inputs.  Old free values at the real anchor are unchanged, and
the actual real formulas and their ordinary derivatives are the same.
Proceed through the finite old recipe.  A raw operation is available by
Theorem~\ref{thm:calculus}.  Every previously positive normalization is
still bounded on its real ordinary neighborhood, so
Lemma~\ref{lem:calculus-positive-bounds} proves its positivity on the
refined flag.  Its old ordinary limit identifies its exponent-zero germ.
Thus a previously regular analytic or implicit step retains its limiting
Jacobian and can be performed by the same theorem.  The additional
passive inputs are held fixed in these comparisons.  Induction proves
the assertion for the entire recipe.  Unused inserted levels contribute
only identity transitions, and later pullback introduces their actual
dependence in the prescribed order.  Preservation of the old free nodes
is supplied by Lemma~\ref{lem:flags-saturation} in the applications.
\end{proof}

The hypothesis about primitive terminal germs in
Lemma~\ref{lem:calculus-refinement} is substantive.  It is verified below
by fixed formal recursions, compatible bounded-charge coefficients, or
canonically normalized infinite-anchor germs.  The chart calculus
preserves that property; it cannot manufacture it from an unknown
exponentially small ambiguity.
